\chapter{Imprimir y escanear}\label{cap:12-imprimir}

Lázaro usa \termino{CUPS}, el sistema de impresión de Linux y de macOS
\cite{cups}. Al instalar revisa las impresoras conectadas y deja puestos
los controladores de las marcas más comunes: HP, Epson, Brother, Samsung y
Xerox.

\section{Agregar la impresora}

\subsection{Por cable USB}
Conéctala y enciéndela. En la mayoría de los casos se agrega sola en
unos segundos. Para comprobarlo, imprime cualquier cosa: la impresora
debe aparecer en la lista.

\subsection{Por wifi o red}
Las impresoras modernas que están en la misma red wifi se detectan solas
y aparecen en la lista al imprimir. Si no aparece:
\begin{pasos}
  \item Abre \menu{Aplicaciones > Sistema > Configuración de impresión}.
  \item Pulsa \menu{Añadir} y escribe tu contraseña.
  \item Elige la impresora en la lista \emph{Impresora de red}. Si no
        está, escribe la dirección IP que muestra la propia impresora en su
        pantalla o en la hoja de configuración que imprime.
  \item Elige la marca y el modelo, y pulsa \menu{Aplicar}.
  \item Imprime una página de prueba.
\end{pasos}

\captura{impresoras}{La ventana de Configuración de impresión}

\begin{consejo}
Si conectas una impresora después de instalar y le falta su controlador,
aparece el aviso \enquote{Falta algo para tu hardware}: pulsa
\menu{Instalar} (sección~\ref{sec:controladores}). ¿Tu modelo funciona en
Linux? La base de datos de OpenPrinting lo dice \cite{openprinting}.
\end{consejo}

\section{Imprimir}

En cualquier programa: \menu{Archivo > Imprimir} o \keys{\ctrl + P}.
Revisa la impresora elegida, la cantidad de copias y las páginas (por
ejemplo, \texttt{1-3,5}). En \emph{Propiedades} o \emph{Configuración de
página} se elige doble cara, blanco y negro o la calidad.

\section{Escanear}

Lázaro instala el programa \emph{Escáner de documentos} (Simple Scan) y
los controladores de escáner de Linux, llamados \termino{SANE}
\cite{sane}. Funciona con la mayoría de los escáneres USB y de las
impresoras multifunción en red.

\begin{pasos}
  \item Abre \menu{Aplicaciones > Gráficos > Escáner de documentos}.
  \item Pon la hoja y pulsa \menu{Escanear}. Elige arriba \emph{Texto} para
        documentos o \emph{Foto} para fotografías.
  \item Para varias hojas, pulsa \menu{Escanear} otra vez: se agregan como
        páginas nuevas.
  \item Pulsa \menu{Guardar}: elige PDF para documentos, JPG para fotos.
\end{pasos}

Las multifunción HP necesitan además su paquete \emph{hplip}, que Lázaro
instala cuando detecta una impresora HP.

\fuentes
